% Part: sets-functions-relations
% Chapter: functions
% Section: basics

\documentclass[../../../include/open-logic-section]{subfiles}

\begin{document}

\olfileid[pt]{sfr}{fun}{bas}
\olsection{Noções Básicas}

\begin{explain}
Uma \emph{função} é um mapa que envia cada !!{element} de um dado conjunto
para um !!{element} específico nalgum (outro) conjunto dado. Por exemplo, a
operação de somar~$1$ define uma função: cada número~$n$ é mapeado
para um único número~$n+1$.

Mais em geral, funções podem tomar pares, triplos, etc., como entradas e
devolver algum tipo de saída. Muitas funções nos são familiares da
aritmética básica. Por exemplo, a adição e a multiplicação são
funções. Tomam dois números e devolvem um terceiro.

Neste sentido matemático e abstrato, uma função é uma \emph{caixa
preta}: só importa que saída fica associada a que entrada, não o
método de calcular a saída.
\end{explain}

\begin{defn}[Função]
Uma \emph{função} $f \colon A \to B$ é um mapeamento de cada !!{element}
de~$A$ para um !!{element} de~$B$.

Chamamos a $A$ o \emph{domínio} de~$f$ e a $B$ o \emph{contradomínio}
de~$f$. Os !!{element}s de~$A$ chamam-se entradas ou \emph{argumentos}
de~$f$, e o !!{element} de~$B$ que fica associado a um argumento~$x$
por~$f$ chama-se o \emph{valor de~$f$} para o argumento~$x$,
escrevendo-se~$f(x)$.

A \emph{imagem} $\ran{f}$ de~$f$ é o subconjunto do contradomínio
constituído pelos valores de~$f$ para algum argumento; $\ran{f} =
\Setabs{f(x)}{x \in A}$.
\end{defn}

O diagrama na \olref{fig:function} pode ajudar a pensar em funções. A elipse
à esquerda representa o \emph{domínio} da função; a elipse à
direita representa o \emph{contradomínio}; e uma seta
aponta de um \emph{argumento} no domínio para o \emph{valor}
correspondente no contradomínio.

\begin{figure}[H]
  \olasset{assets/diagrams/function.tikz}
  \caption{Uma função é um mapeamento de cada !!{element} de um conjunto para
    !!a{element} de outro. Uma seta vai de um argumento no
    domínio para o valor correspondente no contradomínio.}
  \ollabel{fig:function}
\end{figure}

\begin{ex}
A multiplicação toma pares de números naturais como entradas e mapeia-os
para números naturais como saídas, logo vai de $\Nat \times \Nat$ (o
domínio) para $\Nat$ (o contradomínio). Como se vê, a imagem é também
$\Nat$, pois todo $n \in \Nat$ é $n \times 1$.
\end{ex}

\begin{ex}
A multiplicação é uma função porque associa cada entrada---cada par
de números naturais---a uma única saída: $\times \colon \Nat^2 \to
\Nat$. Em contraste, mapear um número natural~$n$ para um real~$x$ tal
que $x^2 = n$ não é funcional, pois cada inteiro positivo~$n$ tem duas
raízes quadradas: $\sqrt{n}$ e~$-\sqrt{n}$. Podemos torná-la funcional
devolvendo só a raiz quadrada positiva: $\sqrt{\phantom{X}} \colon
\Nat \to \Real$.
\end{ex}

\begin{ex}
A relação que associa cada aluno de uma turma à sua nota final
é uma função---nenhum aluno pode ter duas notas finais diferentes na
mesma turma. A relação que associa cada aluno de uma turma aos seus
progenitores não é uma função: os alunos podem ter zero, ou dois, ou mais
progenitores.
\end{ex}

\begin{explain}
Podemos definir funções especificando de modo preciso qual é o
valor da função para cada argumento possível. Formas de o fazer são dar uma fórmula, descrever um método para calcular
o valor, ou listar os valores para cada argumento. Como quer que as funções sejam definidas, temos de garantir que, para cada argumento, especificamos um,
e apenas um, valor.
\end{explain}


\begin{ex}
Seja $f \colon \Nat \to \Nat$ definida por $f(x) = x+1$. Esta
é uma definição que especifica $f$ como uma função que toma
números naturais e devolve números naturais. Diz-nos que, dado um
número natural~$x$, $f$ devolverá o seu sucessor~$x+1$.
Neste caso, o contradomínio $\Nat$ não é a imagem de~$f$, pois o
número natural~$0$ não é sucessor de nenhum número natural. A
imagem de~$f$ é o conjunto de todos os inteiros positivos, $\Int^{+}$.
\end{ex}

\begin{ex}\ollabel{examplefunext}
Seja $g \colon \Nat \to \Nat$ definida por $g(x) = x+2-1$. Isto
nos diz que $g$ é uma função que toma números naturais e
devolve números naturais. Dado um número natural~$n$, $g$ devolverá
o predecessor do sucessor do sucessor de~$n$, O.s.\@,
$n+1$.
\end{ex}

\begin{explain}
Acabamos de considerar duas funções, $f$ e $g$, com
\emph{definições} diferentes. Contudo, são a \emph{mesma função}. Afinal,
para qualquer número natural~$n$, temos $f(n) = n+1 = n+2-1 =
g(n)$. Noutras palavras: as nossas definições de $f$ e~$g$ especificam o
mesmo mapeamento por meio de equações diferentes. Implicitamente, pois,
estamos a apoiar-nos num princípio de extensionalidade para funções,
\[
  \text{se }\forall x\, f(x) = g(x)\text{, então }f = g
\]
desde que $f$ e~$g$ partilhem o mesmo domínio e contradomínio.
\end{explain}

\begin{ex}
Também podemos definir funções por casos. Por exemplo, poderíamos definir
$h \colon \Nat \to \Nat$ por
\[
h(x) =
\begin{cases}
  \frac{x}{2} & \text{se $x$ é par} \\
  \frac{x+1}{2} & \text{se $x$ é ímpar.}
\end{cases}
\]
Como todo número natural é par ou ímpar, a saída desta
função será sempre um número natural. Basta lembrar que, se se
define uma função por casos, cada entrada possível tem de cair em
exatamente um caso. Em alguns casos, isto exigirá uma prova de que os
casos são exaustivos e exclusivos.
\end{ex}

\end{document}
